
To obtain proper estimates for the ESLP, we required:

\begin{enumerate}
    \item Demand: Time series of energy demands for different centres across the UK
    \item Setup Costs: Estimates of UK Renewable Energy Site Setup Costs
    \item Supply Site Capacities: Estimates of UK Renewable Energy Site capacities
    \item Supply Costs: Estimates of generation and transmission costs for different UK Renewable Energy 
\end{enumerate}

Item 1 was obtained from public UK government data on annual UK
electricity consumption statistics across 296 regional authorities from 2005 to 2022 \cite{sub_national_electricity_consumption}.
This provided a timeseries of annual electricity demand across different regional authorities.
Items 2 and 3 were obtained from UK government datasets on all current and completed UK Renewable 
Energy Plants construction projects exceeding 1MW (1880 in total). This included site locations,
technology types e.g. Hydro, Wind, Solar, and capacity.
Based on technology types and capacity, an estimate of setup cost was obtained by combining these
with energy setup capital costs for different technology types, obtained from \cite{nalley2022annual}. Annual capacities 
were obtained by scaling capacity up to cover a year.

\begin{equation}
    \text{Setup Cost}_{USD} = \text{Capacity}_{MWh} \times \text{Capital Cost}_{USD KWh^{-1}} \times 1000 
\end{equation}.

Generation costs in item 3 were obtained from \cite{irena2020}, which provides generation costs (USD) per KWH. 
For transmission, we assumed that all transport was done through overhead wires and obtain an estimate
of 0.0258 USD per km per MWh \cite{desantis2021cost}. A cost per unit of energy between any supply and 
demand locations could then be estimated by calculating the geodesic between two locations. \par

We obtained therefore a list of locations and their energy requirements over time, a list of supply sites,
their setup costs and their capacities, and finally an estimate of supply costs between all demand sites and 
supply sites.
